\section{Stock Profile Parameter Reference}
\label{app:stockparams}

This appendix is the normative parameter reference for the shipping profile
bundle. Every number here is consumed by an equation in the body, and every
equation in the body that requires a stock-specific constant draws it from
here. The tables are the authoritative source; the profile JSON files in the
application bundle are generated from them by the profile-fitting tool, never
hand-edited.

\subsection{Status of These Values}

The values are \emph{provisional fits}. Those marked \textbf{M} in
Table~\ref{tab:provenance} are least-squares fits to published sensitometric
data and are expected to satisfy AC-1 without further work; those marked
\textbf{P} are partial fits in which some parameters are inherited from a
family default because published data is incomplete; those marked \textbf{E}
are engineering estimates awaiting the measurement campaign of
Section~\ref{sec:roadmap}. A stock does not ship until it reaches \textbf{M}
and passes AC-1 (Section~\ref{sec:validation}); this is the mitigation for
risk T-01.

\subsection{Units and Conventions}

\begin{itemize}
\item Densities are Status M diffuse densities for negative stocks and
Status A for print stocks, per ISO 5-3 geometry.
\item Log exposures $x = \logE$ are in $\log_{10}$ lux-seconds at the film
plane, so ISO speed follows from \eqref{eq:speedpoint} rather than being
carried separately. The scene-to-film mapping is \eqref{eq:isoshift}.
\item All spatial parameters are micrometres at the film plane and are
converted to pixels by \eqref{eq:formatscale} at render time.
\item Reversal stocks carry $\gamma < 0$. The characteristic curve
\eqref{eq:charcurve} is unchanged --- $\softp{a}{\cdot}$ is monotone in $u$, and
$u = \gamma(x - x_0)$ is decreasing when $\gamma < 0$, so density decreases with
exposure as a positive image requires. The well-formedness test
\eqref{eq:wellformed} is applied to $|\gamma|$, and the sign of the speed-point
offset in \eqref{eq:speedpoint} inverts accordingly. Reversal stocks default to
the bypass print profile, because a transparency is viewed directly.
\end{itemize}

\subsection{Parameter Inventory}

Table~\ref{tab:paraminventory} lists every stock-carried parameter, the
equation that consumes it, and whether it is per-record or scalar. A profile
that omits a parameter inherits its family default; a profile that carries a
parameter not in this table fails validation at load.

\begin{table}[htbp]
\caption{Stock Profile Parameter Inventory}
\label{tab:paraminventory}
\begin{center}
\renewcommand{\arraystretch}{1.12}
\footnotesize
\begin{tabularx}{\columnwidth}{@{}llc>{\raggedright\arraybackslash}X@{}}
\toprule
\textbf{Symbol} & \textbf{Eq.} & \textbf{Rec.} & \textbf{Quantity} \\
\midrule
$\Dmin$        & \eqref{eq:charcurve}   & \checkmark & Base + fog density \\
$\Delta D$     & \eqref{eq:charcurve}   & \checkmark & Density range \\
$\gamma$       & \eqref{eq:ucoord}      & \checkmark & Straight-line gamma \\
$x_0$          & \eqref{eq:ucoord}      & \checkmark & Extrapolated speed point \\
$\kappa_t$     & \eqref{eq:charcurve}   & \checkmark & Toe softness \\
$\kappa_s$     & \eqref{eq:charcurve}   & \checkmark & Shoulder softness \\
$m$            & \eqref{eq:maskdepletion} & \checkmark & Mask depletion \\
$\Delta x^{\mathrm{(bal)}}$ & \eqref{eq:tungsten} & \checkmark & Layer speed balance \\
$\mathbf{w}_{\mathrm{spectral}}$ & \eqref{eq:pan} & --- & Panchromatic weights \\
$\gamma_\infty/\gamma^{(0)}$ & \eqref{eq:gammatime} & --- & Development ceiling \\
$\phi_0, \beta_f$ & \eqref{eq:fog}      & \checkmark & Fog growth \\
$\alpha_s$     & \eqref{eq:speedshift}  & --- & Speed recovery \\
$\tau_t, \tau_D$ & Sec.~\ref{sec:development} & --- & Toe / range drift \\
$\mathbf{A}$   & \eqref{eq:dirmatrix}   & $3{\times}3$ & DIR coupling \\
$\sigma_1^{\mathrm{DIR}}, \sigma_2^{\mathrm{DIR}}, w_2$ & \eqref{eq:twoscale} & --- & Diffusion scales \\
$\mu$          & \eqref{eq:activityweight} & --- & Activity exponent \\
$G_S$          & \eqref{eq:selwyn}      & \checkmark & Selwyn granularity \\
$\sigma_1^{g}, \chi, \zeta$ & \eqref{eq:grainkernel} & --- & Grain kernel \\
$\nu_1, \nu_2$ & \eqref{eq:grainvar}    & --- & Granularity shape \\
$\varphi$      & \eqref{eq:graincorr}   & --- & Grain chroma correlation \\
$\alpha_h, \ell_R, \omega$ & \eqref{eq:halpsf} & --- & Halation \\
$\boldsymbol{\beta}$ & Tab.~\ref{tab:halparams} & \checkmark & Halation bias \\
$\mathbf{C}$   & \eqref{eq:cmatrix}     & $3{\times}3$ & Printing density \\
$\mathbf{D}'_{\mathrm{aim}}$ & \eqref{eq:aimbalance} & \checkmark & Print aim density \\
\bottomrule
\end{tabularx}
\end{center}
\end{table}

\subsection{Negative Characteristic Curve Parameters}

Table~\ref{tab:negparams} gives the green-record parameters of
\eqref{eq:charcurve} for the ten launch stocks at the manufacturer's normal
process ($A = 1$). The red and blue records are obtained by applying the
per-record offsets of Table~\ref{tab:recordoffsets}; this factoring is used
rather than thirty independent numbers because the offsets are structurally
similar across a stock family and because crossover --- non-parallelism of the
three records --- is then an explicit, reviewable quantity rather than an
accident of three separate fits.

\begin{table*}[htbp]
\caption{Green-Record Characteristic Curve Parameters at Normal Process}
\label{tab:negparams}
\begin{center}
\renewcommand{\arraystretch}{1.15}
\footnotesize
\begin{tabular}{@{}llrrrrrrrl@{}}
\toprule
\textbf{Profile ID} & \textbf{Display name} & \textbf{ISO} & $\gamma_G$ & $\Delta D_G$ & $x_{0,G}$ & $\kappa_{t,G}$ & $\kappa_{s,G}$ & $\Dmin{}_{,G}$ & \textbf{Process} \\
\midrule
\texttt{neg.portra160}  & Portra 160-type       & 160 & $\phantom{-}0.58$ & $1.86$ & $-2.31$ & $0.155$ & $0.115$ & $0.90$ & C-41 \\
\texttt{neg.portra400}  & Portra 400-type       & 400 & $\phantom{-}0.63$ & $1.90$ & $-2.71$ & $0.140$ & $0.110$ & $0.92$ & C-41 \\
\texttt{neg.gold200}    & Gold 200-type         & 200 & $\phantom{-}0.68$ & $1.78$ & $-2.41$ & $0.150$ & $0.098$ & $0.88$ & C-41 \\
\texttt{neg.ektar100}   & Ektar 100-type        & 100 & $\phantom{-}0.72$ & $2.02$ & $-2.11$ & $0.128$ & $0.092$ & $0.86$ & C-41 \\
\texttt{neg.superia400} & Superia 400-type      & 400 & $\phantom{-}0.65$ & $1.84$ & $-2.71$ & $0.145$ & $0.104$ & $0.94$ & C-41 \\
\texttt{neg.v3\_500t}   & Vision3 500T-type     & 500 & $\phantom{-}0.55$ & $2.06$ & $-2.81$ & $0.168$ & $0.130$ & $0.96$ & ECN-2 \\
\texttt{rev.velvia50}   & Velvia 50-type        & 50  & $-1.95$           & $3.10$ & $-1.80$ & $0.090$ & $0.140$ & $0.10$ & E-6 \\
\texttt{rev.provia100}  & Provia 100F-type      & 100 & $-1.72$           & $2.95$ & $-2.10$ & $0.105$ & $0.150$ & $0.09$ & E-6 \\
\texttt{mono.trix400}   & Tri-X 400-type        & 400 & $\phantom{-}0.62$ & $1.92$ & $-2.71$ & $0.160$ & $0.200$ & $0.22$ & B\&W \\
\texttt{mono.hp5}       & HP5 Plus-type         & 400 & $\phantom{-}0.58$ & $1.88$ & $-2.71$ & $0.170$ & $0.215$ & $0.20$ & B\&W \\
\bottomrule
\multicolumn{10}{@{}p{\textwidth}@{}}{\footnotesize Every row satisfies the well-formedness constraint \eqref{eq:wellformed}, $\Delta D \geq 4(\kappa_t + \kappa_s)$, with the smallest margin at \texttt{mono.hp5} ($1.88$ against a required $1.54$) --- monochrome stocks sit closest to the bound because their long shoulders demand a large $\kappa_s$, which is also why the constraint binds first there when a user pushes development. Monochrome stocks carry a single record and the $G$ column is the whole profile; the panchromatic weights of \eqref{eq:pan} are $(0.30, 0.59, 0.11)$ for both, with the orthochromatic variant $(0.02, 0.68, 0.30)$ available as a process option.}
\end{tabular}
\end{center}
\end{table*}

\begin{table}[htbp]
\caption{Per-Record Offsets Relative to the Green Record}
\label{tab:recordoffsets}
\begin{center}
\renewcommand{\arraystretch}{1.15}
\footnotesize
\begin{tabular}{@{}lrrrrrr@{}}
\toprule
& \multicolumn{2}{c}{\textbf{Color negative}} & \multicolumn{2}{c}{\textbf{Transparency}} & \multicolumn{2}{c}{\textbf{Mono}} \\
\cmidrule(lr){2-3}\cmidrule(lr){4-5}\cmidrule(lr){6-7}
\textbf{Offset} & $R$ & $B$ & $R$ & $B$ & $R$ & $B$ \\
\midrule
$\Delta\gamma$      & $-0.02$ & $+0.03$ & $-0.06$ & $+0.08$ & --- & --- \\
$\Delta(\Delta D)$  & $+0.05$ & $-0.05$ & $+0.04$ & $-0.07$ & --- & --- \\
$\Delta x_0$        & $-0.05$ & $+0.05$ & $-0.03$ & $+0.04$ & --- & --- \\
$\Delta\kappa_t$    & $+0.010$ & $-0.010$ & $+0.008$ & $-0.006$ & --- & --- \\
$\Delta\kappa_s$    & $+0.010$ & $\phantom{+}0.000$ & $+0.006$ & $-0.004$ & --- & --- \\
$\Dmin$ (absolute)  & $0.58$ & $1.28$ & $0.11$ & $0.12$ & --- & --- \\
$m$ (mask depl.)    & $0.00$ & $0.10$ & $0.00$ & $0.00$ & --- & --- \\
\bottomrule
\multicolumn{7}{@{}p{0.95\columnwidth}@{}}{\footnotesize The $\Dmin$ row is absolute, not an offset, because it is the orange mask and is the one quantity that must be read directly. The green-record value is in Table~\ref{tab:negparams}; $m_G = 0.06$. Transparency stocks have no mask, so their $\Dmin$ spread is only base tint.}
\end{tabular}
\end{center}
\end{table}

The small $\Delta\gamma$ values are deliberate: three records with identical
gamma have no crossover at all, which is unrealistically clean, while the
$\pm 0.03$ spread here reproduces the mild crossover that makes a real negative
resist a single global balance correction in deep shadow. Transparency stocks
carry roughly three times the spread, which is why slide film is unforgiving of
mixed lighting.

\subsection{Tungsten Balance}

\texttt{neg.v3\_500t} is the only launch stock with a non-zero layer speed
balance in \eqref{eq:tungsten}:

\begin{equation}
\Delta x^{\mathrm{(bal)}}_{\mathrm{v3\_500t}} = (-0.29,\; 0.00,\; +0.42)^\top ,
\end{equation}

derived from the ratio of layer sensitivities integrated against a 3200\,K
Planckian SPD versus D55. Shot under daylight without an 85B correction, this
places the blue record $0.42$ in log exposure higher and the red record $0.29$
lower, which drives blue toward its shoulder in bright areas while red sits in
its toe --- producing a cast that changes character with luminance rather than
a uniform tint. Shipping values for the remaining balance and white-balance
constants are in Table~\ref{tab:balance}.

\subsection{Chemistry and Development}

The chemistry profiles of Table~\ref{tab:chemistry} supply $(\gamma_\infty /
\gamma^{(0)}, \phi_0, \beta_f, \alpha_s, \rho)$. Stocks override only where
their measured push response departs from the family, and carry a per-record
fog coefficient vector because the blue-sensitive layer fogs fastest:

\begin{equation}
\boldsymbol{\phi}_0 / \phi_0 = (0.82,\; 1.00,\; 1.34)
\end{equation}

for all C-41 and ECN-2 stocks, $(0.90, 1.00, 1.15)$ for E-6, and unity for
monochrome. The remaining development constants are family-wide: $\tau_t =
0.30$, $\tau_D = 0.12$, $E_a/R = 8.4\times10^{3}$\,K, $t_0 = 195$\,s at $T_0 =
311.15$\,K (38\,\textdegree C) for C-41, $t_0 = 180$\,s at $314.65$\,K
(41.5\,\textdegree C) for E-6 first development, and $t_0 = 180$\,s at
$314.75$\,K (41.6\,\textdegree C) for ECN-2. Agitation efficiency uses
$\eta_\infty = 1.22$, $\eta_0 = 0.61$, $a_c = 0.85$.

\subsection{Interlayer Inhibition}

The color negative coupling matrix is \eqref{eq:dirmatrix}. Per-stock variation
is carried as a scalar multiplier on $\mathbf{A}$ rather than as ten
independent matrices, because the published data does not support
distinguishing the individual entries and pretending otherwise would be
dishonest about the fit quality:

\begin{center}
\renewcommand{\arraystretch}{1.1}
\footnotesize
\begin{tabular}{@{}lrrr@{}}
\toprule
\textbf{Family} & \textbf{$\mathbf{A}$ scale} & $\sigma_1$ ($\mu$m) & $\sigma_2$ ($\mu$m) \\
\midrule
Portra-type (low contrast, high DIR) & $1.15$ & $1.2$ & $6.0$ \\
Consumer color negative              & $0.85$ & $1.3$ & $6.5$ \\
Saturated color negative             & $1.30$ & $1.1$ & $5.5$ \\
Motion picture negative              & $1.20$ & $1.2$ & $6.2$ \\
Transparency                         & $0.40$ & $1.0$ & $4.5$ \\
Monochrome (scalar $a = 0.34$)       & $1.00$ & $1.4$ & $7.0$ \\
\bottomrule
\end{tabular}
\end{center}

\noindent with $w_2 = 0.38$ and $\mu = 1.0$ throughout. Agitation modulates
$\sigma_2$ and $w_2$ inversely per Section~\ref{sec:interlayer}; at the minimum
agitation setting $\sigma_2$ reaches $1.9\times$ its nominal value and $w_2$
reaches $0.55$, which is the stand-development limit.

\subsection{Grain Parameters}

\begin{table*}[htbp]
\caption{Grain Synthesis Parameters}
\label{tab:grainparams}
\begin{center}
\renewcommand{\arraystretch}{1.15}
\footnotesize
\begin{tabular}{@{}lrrrrrrl@{}}
\toprule
\textbf{Profile ID} & $G_S \times 10^{3}$ & $\sigma_1^{g}$ ($\mu$m) & $\chi$ & $\zeta$ & $\varphi$ & $(\nu_1, \nu_2)$ & \textbf{Mode at 1:1} \\
\midrule
\texttt{neg.portra160}  & $3.0$  & $0.85$ & $0.12$ & $3.0$ & $0.15$ & $(1.00, 1.00)$ & Gaussian \\
\texttt{neg.portra400}  & $4.0$  & $1.05$ & $0.16$ & $3.0$ & $0.15$ & $(1.00, 1.00)$ & Gaussian \\
\texttt{neg.gold200}    & $5.5$  & $1.10$ & $0.20$ & $3.2$ & $0.18$ & $(1.00, 0.92)$ & Gaussian \\
\texttt{neg.ektar100}   & $3.4$  & $0.80$ & $0.10$ & $2.8$ & $0.12$ & $(1.00, 1.00)$ & Gaussian \\
\texttt{neg.superia400} & $4.2$  & $1.02$ & $0.18$ & $3.0$ & $0.16$ & $(1.00, 0.96)$ & Gaussian \\
\texttt{neg.v3\_500t}   & $5.0$  & $1.15$ & $0.17$ & $3.1$ & $0.14$ & $(1.00, 1.00)$ & Gaussian \\
\texttt{rev.velvia50}   & $9.0$  & $0.95$ & $0.14$ & $2.9$ & $0.35$ & $(1.00, 0.80)$ & Gaussian \\
\texttt{rev.provia100}  & $8.0$  & $0.98$ & $0.15$ & $2.9$ & $0.35$ & $(1.00, 0.84)$ & Gaussian \\
\texttt{mono.trix400}   & $17.0$ & $1.55$ & $0.28$ & $3.4$ & $1.00$ & $(1.00, 0.45)$ & Gaussian \\
\texttt{mono.hp5}       & $15.0$ & $1.48$ & $0.26$ & $3.3$ & $1.00$ & $(1.00, 0.48)$ & Gaussian \\
\bottomrule
\multicolumn{8}{@{}p{\textwidth}@{}}{\footnotesize $G_S$ is RMS granularity through a $48\,\mu$m aperture, scaled by $1000$, as published; it pins $\lambda\mathbb{E}[a^2]$ through \eqref{eq:selwyn}, leaving $(\sigma_1^{g}, \chi)$ as the only free parameters, fitted to the published Wiener spectrum by \eqref{eq:wienerclosed}. The mode column is the synthesis mode selected by $\Lambda = \lambda\Delta A_{\mathrm{px}}$ at 135 format and 48\,MP; Poisson mode engages for all stocks at 4$\times$5 format below about 6\,MP, and for the monochrome stocks at Super\,8.}
\end{tabular}
\end{center}
\end{table*}

\subsection{Halation Parameters}

Table~\ref{tab:halparams} gives the family defaults. Per-stock overrides are
confined to three quantities, because the remainder are properties of the base
and antihalation layer rather than of the emulsion:

\begin{center}
\renewcommand{\arraystretch}{1.1}
\footnotesize
\begin{tabular}{@{}lrrrl@{}}
\toprule
\textbf{Profile} & $\alpha_h$ & $\ell_R$ ($\mu$m) & $\omega$ & \textbf{Antihalation} \\
\midrule
\texttt{neg.portra160}  & $0.16$ & $92$  & $0.04$ & Dyed base \\
\texttt{neg.portra400}  & $0.18$ & $95$  & $0.05$ & Dyed base \\
\texttt{neg.gold200}    & $0.21$ & $101$ & $0.07$ & Dyed base \\
\texttt{neg.ektar100}   & $0.14$ & $88$  & $0.04$ & Dyed base \\
\texttt{neg.superia400} & $0.19$ & $97$  & $0.06$ & Dyed base \\
\texttt{neg.v3\_500t}   & $0.05$ & $86$  & $0.01$ & Remjet \\
\texttt{neg.v3\_500t.xr}& $0.55$ & $118$ & $0.22$ & Remjet removed \\
\texttt{rev.velvia50}   & $0.10$ & $80$  & $0.03$ & Dyed base \\
\texttt{rev.provia100}  & $0.11$ & $82$  & $0.03$ & Dyed base \\
\texttt{mono.trix400}   & $0.12$ & $90$  & $0.05$ & Dyed base \\
\texttt{mono.hp5}       & $0.13$ & $93$  & $0.06$ & Dyed base \\
\bottomrule
\end{tabular}
\end{center}

\noindent The \texttt{.xr} variant is the motion picture stock with its remjet
backing removed for still-camera use, which is the reason that particular look
--- a broad red halo around every specular highlight --- became recognizable.
It is a separate profile rather than a slider, because it is a separate
physical product. Monochrome stocks use $\boldsymbol{\beta} = (1,1,1)$; all
others use the values in Table~\ref{tab:halparams}.

\subsection{Print Stock Parameters}

\begin{table}[htbp]
\caption{Print Stock Curve Parameters}
\label{tab:printparams}
\begin{center}
\renewcommand{\arraystretch}{1.15}
\footnotesize
\begin{tabular}{@{}lrrrrr@{}}
\toprule
\textbf{Profile ID} & $\gamma'$ & $\Delta D'$ & $\kappa'_t$ & $\kappa'_s$ & $\Dmin'$ \\
\midrule
\texttt{prt.2383}   & $2.80$ & $2.42$ & $0.220$ & $0.060$ & $0.06$ \\
\texttt{prt.2393}   & $3.05$ & $2.58$ & $0.185$ & $0.052$ & $0.05$ \\
\texttt{prt.3513}   & $2.72$ & $2.35$ & $0.245$ & $0.068$ & $0.07$ \\
\texttt{prt.3521}   & $2.90$ & $2.48$ & $0.205$ & $0.058$ & $0.06$ \\
\texttt{prt.bypass} & \multicolumn{5}{l}{No transfer; \eqref{eq:printcurve} skipped, normalize only} \\
\bottomrule
\end{tabular}
\end{center}
\end{table}

The aim density triple of \eqref{eq:aimbalance} is
$\mathbf{D}'_{\mathrm{aim}} = (1.09,\; 1.06,\; 1.03)^\top$ for all four print
stocks, following laboratory aim density practice; the small red-to-blue
gradient is the standard allowance for projector lamp color temperature and is
what makes the neutral axis of a printed frame lean very slightly warm.

Printing density matrices are the base $\mathbf{C}$ of \eqref{eq:cmatrix} with
the off-diagonal terms scaled per stock: $\times 1.00$ for \texttt{prt.2383},
$\times 1.18$ for \texttt{prt.2393}, $\times 0.92$ for \texttt{prt.3513}, and
$\times 1.09$ for \texttt{prt.3521}, with the Fuji-type stocks additionally
carrying a $+0.014$ adjustment on $C_{GB}$ that is the numerical form of their
distinct green rendering. Silver retention uses $\Delta D'_{\mathrm{Ag}} =
0.90$ with $\varrho$ defaulting to $0$; the ENR preset is $\varrho = 0.45$ and
full bleach bypass is $\varrho = 1.0$.

\subsection{Profile File Format}

Profiles are versioned JSON, one file per stock, loaded into
\texttt{stock\_profile.params\_json} (Appendix~\ref{app:ddl}) at first launch
and thereafter from the database. Unspecified fields inherit the family
default named by \texttt{inherits}.

\begin{lstlisting}[language=Swift, caption={Stock profile file (abridged).}, label={lst:profilejson}]
{
  "id": "neg.portra400",
  "profileVersion": 3,
  "kind": "negative",
  "displayName": "Portra 400-type",
  "inherits": "family.colorneg.portra",
  "fitStatus": "M",
  "iso": 400,
  "curve": {
    "green": { "gamma": 0.63, "deltaD": 1.90, "x0": -2.71,
               "kappaT": 0.140, "kappaS": 0.110, "dMin": 0.92 },
    "offsets": {
      "red":  { "gamma": -0.02, "deltaD":  0.05, "x0": -0.05,
                "kappaT": 0.010, "kappaS": 0.010, "dMin": 0.58 },
      "blue": { "gamma":  0.03, "deltaD": -0.05, "x0":  0.05,
                "kappaT": -0.010, "kappaS": 0.000, "dMin": 1.28 }
    },
    "maskDepletion": [0.00, 0.06, 0.10],
    "balanceShift":  [0.00, 0.00, 0.00]
  },
  "chemistry": { "default": "chem.c41", "fogPerRecord": [0.82, 1.00, 1.34] },
  "interlayer": { "matrixScale": 1.15, "sigma1um": 1.2, "sigma2um": 6.0,
                  "w2": 0.38, "mu": 1.0 },
  "grain": { "selwyn": 4.0, "sigma1um": 1.05, "chi": 0.16, "zeta": 3.0,
             "phi": 0.15, "nu": [1.00, 1.00] },
  "halation": { "alpha": 0.18, "lengthRedUm": 95.0, "omega": 0.05,
                "bias": [1.00, 0.42, 0.22], "antihalation": "dyedBase" },
  "defaultPrint": "prt.2383",
  "provenance": {
    "source": "manufacturer datasheet, D-log E family + Wiener spectrum",
    "fittedAt": "2026-02-11", "rmsDensityError": 0.021
  }
}
\end{lstlisting}

\subsection{Fit Provenance}

\begin{table}[htbp]
\caption{Fit Status and Data Provenance}
\label{tab:provenance}
\begin{center}
\renewcommand{\arraystretch}{1.15}
\footnotesize
\begin{tabularx}{\columnwidth}{@{}llc>{\raggedright\arraybackslash}X@{}}
\toprule
\textbf{Profile} & \textbf{Fit} & \textbf{RMS $D$} & \textbf{Basis} \\
\midrule
\texttt{neg.portra400}  & M & $0.021$ & Full $D$--$\logE$ family, Wiener spectrum, MTF \\
\texttt{neg.portra160}  & M & $0.019$ & Full family, Wiener spectrum \\
\texttt{neg.ektar100}   & M & $0.024$ & Full family, Wiener spectrum \\
\texttt{neg.v3\_500t}   & M & $0.023$ & Full family, spectral sensitivities \\
\texttt{neg.superia400} & P & $0.038$ & Curves at normal process only \\
\texttt{neg.gold200}    & P & $0.041$ & Curves only; grain estimated \\
\texttt{rev.velvia50}   & M & $0.028$ & Full family \\
\texttt{rev.provia100}  & P & $0.035$ & Curves only \\
\texttt{mono.trix400}   & M & $0.017$ & Full family, published CI vs.\ time \\
\texttt{mono.hp5}       & M & $0.018$ & Full family, published CI vs.\ time \\
\bottomrule
\multicolumn{4}{@{}p{0.95\columnwidth}@{}}{\footnotesize AC-1 requires RMS density error $\leq 0.03$ and maximum error $\leq 0.06$. The three \textbf{P} profiles are the ones that must reach \textbf{M} through the measurement campaign before launch; they are the concrete content of risk T-01, and the fallback is to ship eight validated stocks rather than ten unvalidated ones.}
\end{tabularx}
\end{center}
\end{table}
